\begin{frame}{``우연히 좋은'' 패턴 5가지 (2/2): 4·5}
  \scriptsize
  \begin{tabular}{@{}p{0.16\textwidth}p{0.24\textwidth}p{0.26\textwidth}p{0.30\textwidth}@{}}
    \toprule
    패턴 & 겉보기 & 실제로는 & 반박 가능한 질문(근거) \\
    \midrule
    \textbf{4. 하이퍼파라미터 ``적당히'' 하나} & ``$\alpha=0.5$로 학습'' &
      \textbf{민감성 미검증} --- Ch\,7.2 n 튜닝처럼 $\alpha$가 크면
      (0.5) Q값이 노이즈에 크게 흔들려 정책이 달라질 수 있음 &
      ``$\alpha$를 0.1로 줄이면 \textbf{결론이 변하나요}? 민감성을
      확인했나요?''(Ch\,7.2 튜닝) \\
    \textbf{5. 베이스라인 없이 ``잘 배웠다''} & ``$-13$이 나왔다'' &
      기준점 없이 ``나다''를 말할 수 없음 --- 무작위 정책 평균 리턴
      (약 $-5000$)을 재야 ``실제로 이득인가'' 판정 가능 & ``무작위
      정책의 평균 리턴은? 배운 정책이 그 대비 얼마나 \textbf{나은지
      측정}했나요?''(8.1 M2) \\
    \bottomrule
  \end{tabular}
  \vskip0.7em
  \begin{block}{패턴 2가 특히 RL \emph{특유}인 이유}
    \textbf{비대칭성} --- 탐험 혼입 시 두 알고리즘이 \emph{다른 방향의
    크기}로 망가지는 것. Ch\,6.3 on/off-policy 경로 차이를
    \textbf{정량적}으로 이해해야 짚어낼 수 있는, 이번 학기 개념이
    직접 작동하는 지점.
  \end{block}
\end{frame}
